\subsection{截断积分与主值}
\label{sec:ch5-truncated-integrals-principal-value}

定理~\ref{thm:ch5-calderon-zygmund} 通过假设 $\widehat K\in L^\infty$ 来保证算子在 $L^2$ 上有界. 本节给出直接施加于 $K$ 的条件, 使相应算子在 $L^p$, $1<p<\infty$, 上有界. 对 $K\in L^1_{\mathrm{loc}}(\mathbb R^n\setminus\{0\})$, 定义 $K_{\varepsilon,R}(x)=K(x)\chi_{\{\varepsilon<|x|<R\}}$.

\begin{Proposition}
\label{prop:ch5-truncated-kernel-fourier-bound}
设 $K\in L^1_{\mathrm{loc}}(\mathbb R^n\setminus\{0\})$ 满足
\begin{equation}
\label{eq:ch5-annular-cancellation}
 \left|\int_{a<|x|<b}K(x)\,dx\right|\le A,
 \qquad 0<a<b<\infty,
\end{equation}
\begin{equation}
\label{eq:ch5-annular-size}
 \int_{a<|x|<2a}|K(x)|\,dx\le B,
 \qquad a>0,
\end{equation}
\begin{equation}
\label{eq:ch5-kernel-difference-condition}
 \int_{|x|>2|y|}|K(x-y)-K(x)|\,dx\le C,
 \qquad y\in\mathbb R^n.
\end{equation}
则对所有 $\xi\in\mathbb R^n$, $|\widehat K_{\varepsilon,R}(\xi)|\le C'$, 其中 $C'$ 与 $\varepsilon,R$ 无关.
\end{Proposition}

条件 \eqref{eq:ch5-annular-size} 等价于
\begin{equation}
\label{eq:ch5-first-moment-size}
 \int_{|x|<a}|x|\,|K(x)|\,dx\le B'a,
 \qquad a>0.
\end{equation}
等价性由二进环带分解直接得到.

\begin{Proof}
固定 $\xi$. 当 $\varepsilon<|\xi|^{-1}<R$ 时, 在半径 $|\xi|^{-1}$ 处分割 Fourier 积分为 $I_1+I_2$. 对 $I_1$, 加减 $1$ 后由 \eqref{eq:ch5-annular-cancellation} 和 \eqref{eq:ch5-first-moment-size} 得 $|I_1|\le C(A+B)$. 对 $I_2$, 取 $z=\xi/(2|\xi|^2)$, 则 $e^{2\pi iz\cdot\xi}=-1$. 将 $I_2$ 平移 $z$, 与原积分相加, 得到一个核差积分以及内外边界上的两个薄环带积分. 第一项由 \eqref{eq:ch5-kernel-difference-condition} 控制, 后两项由 \eqref{eq:ch5-annular-size} 控制. 其他两种 $|\xi|^{-1}$ 相对于 $\varepsilon,R$ 的位置只需保留其中一个积分, 论证相同.
\end{Proof}

\hypertarget{chapter-05-page-107}{}
\begin{Corollary}
\label{cor:ch5-uniform-truncation-bounds}
若 $K$ 满足命题~\ref{prop:ch5-truncated-kernel-fourier-bound} 的假设, 则
\[
 \|K_{\varepsilon,R}*f\|_p\le C_p\|f\|_p,
 \qquad 1<p<\infty,
\]
并且
\[
 |\{x\in\mathbb R^n:|K_{\varepsilon,R}*f(x)|>\lambda\}|
 \le\frac{C_1}{\lambda}\|f\|_1,
\]
其中常数与 $\varepsilon,R$ 无关.
\end{Corollary}

$p=2$ 时结论直接来自命题~\ref{prop:ch5-truncated-kernel-fourier-bound}. 其余情形只需注意 \eqref{eq:ch5-kernel-difference-condition} 蕴含一致的 Hörmander 条件
\[
 \int_{|x|>2|y|}|K_{\varepsilon,R}(x-y)-K_{\varepsilon,R}(x)|\,dx\le C.
\]
细节留给读者.

\hypertarget{chapter-05-page-108}{}
当 $K(x)=\Omega(x')/|x|^n$ 齐次时, 条件 \eqref{eq:ch5-annular-cancellation} 和 \eqref{eq:ch5-annular-size} 等价于 $\Omega\in L^1(S^{n-1})$ 且 $\int_{S^{n-1}}\Omega\,d\sigma=0$.

由于 $K_{\varepsilon,R}\in L^1$, $K_{\varepsilon,R}*f$ 对 $f\in L^p$ 有定义. 理想情形是用逐点极限
\[
 Tf(x)=\lim_{\substack{\varepsilon\to0\\R\to\infty}}
 K_{\varepsilon,R}*f(x)
\]
定义奇异积分, 但即便 $f\in\mathcal S$, 这个极限也未必存在. 当 $R\to\infty$ 时已有收敛, 所以问题归结为判断主值分布
\[
 \operatorname{p.v.}K(\varphi)
 =\lim_{\varepsilon\to0}\int_{|x|>\varepsilon}K(x)\varphi(x)\,dx
\]
何时存在.

\begin{Proposition}
\label{prop:ch5-principal-value-existence}
若 $K$ 满足 \eqref{eq:ch5-annular-size}, 则缓增分布 $\operatorname{p.v.}K$ 存在, 当且仅当极限
\[
 \lim_{\varepsilon\to0}\int_{\varepsilon<|x|<1}K(x)\,dx
\]
存在.
\end{Proposition}

\begin{Proof}
若主值分布存在, 取在 $B(0,1)$ 上恒等于 $1$ 的 $\varphi\in\mathcal S$. 分割 $|x|<1$ 与 $|x|>1$ 后, 外部积分由 \eqref{eq:ch5-annular-size} 和 Schwartz 衰减收敛, 故内部截断积分必有极限.

反之, 设该极限为 $L$. 则定义
\[
 \operatorname{p.v.}K(\varphi)=\varphi(0)L
 +\int_{|x|<1}K(x)[\varphi(x)-\varphi(0)]\,dx
 +\int_{|x|>1}K(x)\varphi(x)\,dx.
\]
末项由环带分解收敛. 对中间项使用 $|\varphi(x)-\varphi(0)|\le|x|\|\nabla\varphi\|_\infty$ 及 \eqref{eq:ch5-first-moment-size}, 也得到绝对收敛.
\end{Proof}

\hypertarget{chapter-05-page-109}{}
\begin{Corollary}
\label{cor:ch5-principal-value-operator-bounds}
若在命题~\ref{prop:ch5-truncated-kernel-fourier-bound} 的假设之外再假定
\[
 \lim_{\varepsilon\to0}\int_{\varepsilon<|x|<1}K(x)\,dx
\]
存在, 则
\[
 Tf(x)=\lim_{\varepsilon\to0}\int_{|y|>\varepsilon}K(y)f(x-y)\,dy
\]
可延拓为 $L^p$, $1<p<\infty$, 上的有界算子, 并且是弱型 $(1,1)$ 的.
\end{Corollary}

\begin{Example}[$K(x)=|x|^{-n-it}$ 的应用]
\label{ex:ch5-complex-power-kernel}
函数 $|x|^{-n-it}$ 在 $\mathbb R^n\setminus\{0\}$ 上局部可积, 并满足命题~\ref{prop:ch5-truncated-kernel-fourier-bound} 的假设. 事实上,
\[
 \left|\int_{a<|x|<b}\frac{dx}{|x|^{n+it}}\right|
 =|S^{n-1}|\left|\frac{b^{-it}-a^{-it}}{-it}\right|
 \le\frac2{|t|}|S^{n-1}|,
\]
\[
 \int_{a<|x|<2a}\frac{dx}{|x|^n}=|S^{n-1}|\log2,
 \qquad
 |\nabla K(x)|\le\frac{|n+it|}{|x|^{n+1}}.
\]
因此 $\|K_{\varepsilon,R}*f\|_p\le C_p\|f\|_p$, 且 $C_p$ 与 $\varepsilon,R$ 无关. 但截断积分的极限几乎处处不存在, 因为
\[
 \int_{\varepsilon<|x|<1}\frac{dx}{|x|^{n+it}}
 =|S^{n-1}|\frac{1-\varepsilon^{-it}}{-it},
\]
而 $\varepsilon^{-it}$ 在 $\varepsilon\to0$ 时没有极限. 选取 $\varepsilon_k=e^{-2\pi k/t}$ 则极限存在, 并定义一个强型 $(p,p)$, $1<p<\infty$, 且弱型 $(1,1)$ 的卷积算子. 称相应分布为 $\operatorname{p.v.}|x|^{-n-it}$, 尽管它依赖于序列的选择:
\[
 \operatorname{p.v.}\frac1{|x|^{n+it}}(\varphi)
 =\int_{|x|<1}\frac{\varphi(x)-\varphi(0)}{|x|^{n+it}}\,dx
 +\int_{|x|>1}\frac{\varphi(x)}{|x|^{n+it}}\,dx.
\]
\end{Example}

\hypertarget{chapter-05-page-110}{}
尽管有这样的记号, 这个分布不是 $-n-it$ 次齐次的. 对 $\operatorname{Re}z<n$ 考察局部可积齐次分布 $|x|^{-z}$. 将它写为
\[
 \frac1{|x|^z}(\varphi)
 =\int_{|x|<1}\frac{\varphi(x)-\varphi(0)}{|x|^z}\,dx
 +\int_{|x|>1}\frac{\varphi(x)}{|x|^z}\,dx
 +\frac{|S^{n-1}|}{n-z}\varphi(0),
\]
右端除 $z=n$ 外可延拓到 $\operatorname{Re}z<n+1$. 令 $z=n+it$, 它与上述主值分布相差 Dirac 测度的常数倍. 新分布是唯一的 $-n-it$ 次齐次分布, 且在原点外与 $|x|^{-n-it}$ 一致. 利用 \eqref{eq:ch4-fourier-power-kernel}, 可得
\[
 \left(\operatorname{p.v.}\frac1{|x|^{n+it}}\right)^{\!\wedge}(\xi)
 =\pi^{n/2+it}\frac{\Gamma(-it/2)}{\Gamma((n+it)/2)}
 |\xi|^{it}+\frac1{it}|S^{n-1}|,
\]
其中最后的常数项反映了所选主值与齐次延拓之间的 Dirac 项.
